% GENERATED FILE. Edit canonical lessons, then run npm run generate:books.
% canonical-source-hash: fnv1a64:79848e2ba348a310
% canonical-lessons: IT-C323-c-era-una-volta, IT-C323-racconto, IT-C323-favola, IT-C323-narrare, IT-C323-narratore

\chapter{C'era una volta}
\label{ch:it-c-era-una-volta}
\hlchaptermodality{323}

\begin{chapteropening}
\textbf{By the end of this chapter:} \emph{I can begin and tell a story: once upon a time, the tale, the fable, narrating and the narrator.}

Complete the last lesson of chapter 323: I can begin and tell a story: once upon a time, the tale, the fable, narrating and the narrator.
\end{chapteropening}

\section[c'era una volta]{c'era una volta — once upon a time (there was)}

\label{lesson:IT-C323-c-era-una-volta}



\begin{quote}
Before the new one: say the Italian for the lantern, then the Italian for to dance.
\end{quote}

\subsection*{You'll want to know: c'era una volta}
\textbf{c'era una volta} — \textquotedblleft{}once upon a time (there was)\textquotedblright{}.

\begin{quote}
C'era una volta una ragazza che viveva nel bosco.
\end{quote}

\begin{quote}
\emph{Once upon a time there was a girl who lived in the woods.}
\end{quote}

The opening of almost every fairy tale; \emph{c'era} is the \emph{imperfetto} of \emph{c'è}.

\begin{cousinweb}
\emph{volta} \textquotedblleft{}a time, a turn\textquotedblright{} is built on \emph{voltare}, usually traced to Vulgar Latin \emph{volvitāre}, from Latin \emph{volvere} \textquotedblleft{}to roll, to turn\textquotedblright{} (Lewis \& Short).
\end{cousinweb}

\begin{grammarlens}[title={What you've built}]
One more word for telling, explaining and describing at B1.
\end{grammarlens}

\subsection*{Your turn}
\begin{itemize}
\raggedright
  \item \emph{Say it:} \emph{c'era una volta}
  \item \emph{Say it:} \emph{c'era una volta}, once more
  \item \emph{Say it:} a sentence of your own with \emph{c'era una volta}
  \item \emph{Recall:} say the Italian for the lantern, then the Italian for to dance, then say \emph{c'era una volta} again
\end{itemize}

\begin{tcolorbox}[breakable,colback=teal!4,colframe=teal!35!black,title={Before you move on}]
What does c'era una volta mean? (\textquotedblleft{}Once upon a time (there was)\textquotedblright{}.) Say it once more.
\end{tcolorbox}
\section[il racconto]{il racconto — the story, the short story}

\label{lesson:IT-C323-racconto}



\begin{quote}
Before the new one: say the Italian for to dance, then the Italian for once upon a time.
\end{quote}

\subsection*{You'll want to know: il racconto}
\textbf{il racconto} — \textquotedblleft{}the story, the short story\textquotedblright{}.

\begin{quote}
Il nonno ci leggeva un racconto ogni sera.
\end{quote}

\begin{quote}
\emph{Grandfather used to read us a story every evening.}
\end{quote}

Built on \emph{raccontare}, ‘to tell’.

\begin{cousinweb}
\emph{racconto} is built on \emph{raccontare} \textquotedblleft{}to tell\textquotedblright{}, from \emph{ri-}, \emph{a-} and \emph{contare} \textquotedblleft{}to count, to tell\textquotedblright{}, from Latin \emph{computāre} (Lewis \& Short).
\end{cousinweb}

\begin{grammarlens}[title={What you've built}]
One more word for telling, explaining and describing at B1.
\end{grammarlens}

\subsection*{Your turn}
\begin{itemize}
\raggedright
  \item \emph{Say it:} \emph{il racconto}
  \item \emph{Say it:} \emph{il racconto}, once more
  \item \emph{Say it:} a sentence of your own with \emph{il racconto}
  \item \emph{Recall:} say the Italian for to dance, then the Italian for once upon a time, then say \emph{il racconto} again
\end{itemize}

\begin{tcolorbox}[breakable,colback=teal!4,colframe=teal!35!black,title={Before you move on}]
What does il racconto mean? (\textquotedblleft{}The story, the short story\textquotedblright{}.) Say it once more.
\end{tcolorbox}
\section[la favola]{la favola — the fable, the tale}

\label{lesson:IT-C323-favola}



\begin{quote}
Before the new one: say the Italian for once upon a time, then the Italian for the story.
\end{quote}

\subsection*{You'll want to know: la favola}
\textbf{la favola} — \textquotedblleft{}the fable, the tale\textquotedblright{}.

\begin{quote}
Nella favola la volpe parla con il lupo.
\end{quote}

\begin{quote}
\emph{In the fable the fox talks to the wolf.}
\end{quote}

\emph{Una vita da favola} — a fairy-tale life.

\begin{cousinweb}
\emph{favola} is inherited from Latin \emph{fābula} \textquotedblleft{}a story, a tale\textquotedblright{} (EtymDB), from \emph{fārī} \textquotedblleft{}to speak\textquotedblright{}.
\end{cousinweb}

\begin{grammarlens}[title={What you've built}]
One more word for telling, explaining and describing at B1.
\end{grammarlens}

\subsection*{Your turn}
\begin{itemize}
\raggedright
  \item \emph{Say it:} \emph{la favola}
  \item \emph{Say it:} \emph{la favola}, once more
  \item \emph{Say it:} a sentence of your own with \emph{la favola}
  \item \emph{Recall:} say the Italian for once upon a time, then the Italian for the story, then say \emph{la favola} again
\end{itemize}

\begin{tcolorbox}[breakable,colback=teal!4,colframe=teal!35!black,title={Before you move on}]
What does la favola mean? (\textquotedblleft{}The fable, the tale\textquotedblright{}.) Say it once more.
\end{tcolorbox}
\section[narrare]{narrare — to narrate, to tell (a story)}

\label{lesson:IT-C323-narrare}



\begin{quote}
Before the new one: say the Italian for the story, then the Italian for the fable.
\end{quote}

\subsection*{You'll want to know: narrare}
\textbf{narrare} — \textquotedblleft{}to narrate, to tell (a story)\textquotedblright{}.

\begin{quote}
Il libro narra la storia di un ragazzo.
\end{quote}

\begin{quote}
\emph{The book tells the story of a boy.}
\end{quote}

A literary \emph{raccontare}.

\begin{cousinweb}
\emph{narrare} is a learned borrowing from Latin \emph{nārrāre} \textquotedblleft{}to tell, to relate\textquotedblright{} (Lewis \& Short), built on \emph{gnārus} \textquotedblleft{}knowing\textquotedblright{}.
\end{cousinweb}

\begin{grammarlens}[title={What you've built}]
One more word for telling, explaining and describing at B1.
\end{grammarlens}

\subsection*{Your turn}
\begin{itemize}
\raggedright
  \item \emph{Say it:} \emph{narrare}
  \item \emph{Say it:} \emph{narrare}, once more
  \item \emph{Say it:} a sentence of your own with \emph{narrare}
  \item \emph{Recall:} say the Italian for the story, then the Italian for the fable, then say \emph{narrare} again
\end{itemize}

\begin{tcolorbox}[breakable,colback=teal!4,colframe=teal!35!black,title={Before you move on}]
What does narrare mean? (\textquotedblleft{}To narrate, to tell (a story)\textquotedblright{}.) Say it once more.
\end{tcolorbox}
\section[il narratore]{il narratore — the narrator, the storyteller}

\label{lesson:IT-C323-narratore}



\begin{quote}
Before the new one: say the Italian for the fable, then the Italian for to narrate.
\end{quote}

\subsection*{You'll want to know: il narratore}
\textbf{il narratore} — \textquotedblleft{}the narrator, the storyteller\textquotedblright{}.

\begin{quote}
Il narratore parla in prima persona.
\end{quote}

\begin{quote}
\emph{The narrator speaks in the first person.}
\end{quote}

A woman is \emph{la narratrice}.

\begin{cousinweb}
\emph{narratore} is a learned borrowing from Latin \emph{nārrātor} \textquotedblleft{}a storyteller\textquotedblright{} (Lewis \& Short).
\end{cousinweb}

\begin{grammarlens}[title={What you've built}]
One more word for telling, explaining and describing at B1.
\end{grammarlens}

\subsection*{Your turn}
\begin{itemize}
\raggedright
  \item \emph{Say it:} \emph{il narratore}
  \item \emph{Say it:} \emph{il narratore}, once more
  \item \emph{Say it:} a sentence of your own with \emph{il narratore}
  \item \emph{Recall:} say the Italian for the fable, then the Italian for to narrate, then say \emph{il narratore} again
\end{itemize}

\begin{tcolorbox}[breakable,colback=teal!4,colframe=teal!35!black,title={Before you move on}]
What does il narratore mean? (\textquotedblleft{}The narrator, the storyteller\textquotedblright{}.) Say it once more.
\end{tcolorbox}
